Como fue mencionado en la sección anterior, si bien existe software que permite realizar estas simulaciones, no hay actualmente herramientas libres basadas en el uso de lenguajes y técnicas modernas de modelado que integren las funcionalidades requeridas. \\

Con el objetivo de desarrollar un sistema que se acerque a cumplir estos requerimientos, se va a elaborar un proyecto en tres partes:
\begin{itemize}
    \item Una biblioteca en C++: Encargada de implementar una interfaz de comunicación de baja latencia con el hardware comercial utilizado en la simulación de vehículos (tales como servomotores para la respuesta del volante o Force Feedback, pedales y caja de cambios).
    \item Un modelo computacional en Modelica: Diseñado para simular la física y dinámica del vehículo. Este sistema procesará las entradas del conductor, calculará las fuerzas, momentos y transferencias de carga que afectan al vehículo, y emitirá los resultados de manera continua.
    \item Un entorno visualizador en Unity: Responsable de recibir la información de posicionamiento desde Modelica para renderizarla gráficamente, proporcionando al usuario una retroalimentación visual fluida.
\end{itemize}
Este sistema va a tener dos puntos de desarrollo claves: En primer lugar, la ejecución de la simulación en tiempo real estricto, lo cual es importante para que el modelo sea verdaderamente interactivo y útil para un conductor humano. En segundo lugar, su arquitectura modular. Al utilizar Modelica, un lenguaje de simulación orientado a objetos, resultará sencillo intercambiar diferentes partes y subsistemas del vehículo. Esto permitirá evaluar diversas geometrías de suspensión, modificar la rigidez de resortes y amortiguadores, o alterar componentes enteros con un esfuerzo mínimo. \\
Dada la complejidad de simular un vehículo completo, el desarrollo se concentrará principalmente en la simulación detallada del chasis y la cinemática de la suspensión. Sin embargo, la propia naturaleza de Modelica permitirá escalar el proyecto, dejando abierta la posibilidad de cambiar modelos simplificados (como los de neumáticos, motor o transmisión) por variantes más complejas si los requerimientos lo exigen. Finalmente, el diseño desacoplado asegura que las tres partes del sistema sean intercambiables; por ejemplo, se podría sustituir el visualizador en Unity por un software comercial alternativo, o adaptar la interfaz de los controles a hardware específico, sin comprometer la integridad del núcleo de simulación.